\documentclass[../../main.tex]{subfiles}% 注意这里的文件路径不能用 ./main.tex ,否则用latexmk编译子文件会报错
\graphicspath{{\subfix{./image/}}{\subfix{../../image/}}} % 指定图片引用路径,后续可以直接使用图片文件名
% 如果图片存放在上级目录,则使用 ../../image/ 作为备用路径

% 例如:
% \begin{figure}[H]
% \centering
% \includegraphics[width=0.3\textwidth]{图.png}
% \caption{}
% \label{figure:图}
% \end{figure}
% 注意:上述\label{}一定要放在\caption{}之后,否则引用图片序号会只会显示??.

\begin{document}

\section{存在性与正则性结果}\label{section:ch07-存在性}

解的存在性通过一个逼近程序证明。考虑逼近问题
\begin{align*}
\begin{cases}
\partial_t u_n+\partial_x[\phi(u_n)]=0 & \text{在 }S,\\
u_n=u_{0n} & \text{在 }\mathbb{R}\times\{0\}
\end{cases}
\qquad(Pn)
\end{align*}
研究序列 $\{u_n\}$ 的极限点，我们将证明下列结果。

\begin{theorem}\label{theorem:ch07-7.3.1}
设 (A0)--(A1) 成立。则问题 (P) 有解 $u$，它是问题 (Pn) 的熵解序列 $\{u_n\}$ 的一个极限点。此外，$u$ 是 (P) 的 Young 测度熵解。
\end{theorem}

\begin{theorem}\label{theorem:ch07-7.3.2}
设 (A0)--(A2) 成立，$u$ 为由\cref{theorem:ch07-7.3.1}给出的 (P) 的解。则 $u$ 是 (P) 的熵解。
\end{theorem}

下列定义是\cref{definition:ch06-构造解}的对应物。

\begin{definition}\label{definition:ch07-构造解}
由\cref{theorem:ch07-7.3.1}与\cref{theorem:ch07-7.3.2}给出的 (P) 的解分别称为所构造解与所构造熵解。
\end{definition}

如\cref{section:ch06-存在正则}，(P) 的所构造熵解有若干重要性质。特别地，在附加条件 (A2')（它反过来蕴含 (A1) 与 (A2)）下，它们满足相容条件。这就是下列定理的内容。

\begin{theorem}\label{theorem:ch07-7.3.3}
设 (A0) 与 (A2') 成立，$u$ 为问题 (P) 的所构造熵解。则对每个 $\tau\in(0,T)$：
\begin{enumerate}[(i)]
\item 若 $\phi$ 有界，则 $u$ 在 $[0,\tau]$ 上于每个点 $x_j\in\operatorname{supp}u_s(\cdot,\tau)$ 处满足相容条件 (C$^\pm$)；
\item 若 $\phi$ 无界，则 $u=u_r\in L^\infty(\mathbb{R}\times(\tau,T))$。
\end{enumerate}
\end{theorem}

\begin{remark}\label{remark:ch07-7.3.1}
若 $\phi$ 无界且满足 (A2')，则由文献~\cite{ref67} 定理 1.1 与\cref{theorem:ch07-7.3.3}(ii)，对每个 $u_0\in\mathscr{R}_f^+(\mathbb{R})$ 存在问题 (P) 的唯一熵解，其等待时间 $t_0=0$。事实上，由\cref{theorem:ch07-7.3.3}(ii)，每个所构造熵解 $u$ 是文献~\cite{ref67}意义下的解，因对每个 $\tau\in(0,T)$ 有 $u=u_r\in L^\infty(\mathbb{R}\times(\tau,T))$，且 $\operatorname*{ess\,lim}_{t\to0^+}u(\cdot,t)=u_0$ 在 $\mathscr{R}_f(\mathbb{R})$ 中窄收敛（见\cref{definition:ch05-窄收敛}）。这由\eqref{align:ch07-7.12}与\cref{proposition:ch07-7.1.2}(ii) 得到（另见文献~\cite{ref53} 命题 2，第 38 页）。

当 $\phi$ 有界且 $u_0$ 如 (A0') 时，由文献~\cite{ref17} 命题 3.20，(P) 的每个所构造熵解满足 $u\in C([0,T];\mathscr{R}_f^+(\Omega))$ 且 $u(\cdot,0)=u_0$。故由\cref{theorem:ch07-7.3.3}与\cref{theorem:ch07-唯一性}得下列存在唯一性结果。
\end{remark}

\begin{theorem}\label{theorem:ch07-存在唯一}
设 $\phi$ 有界，(A0')--(A2') 成立。则问题 (P) 存在唯一熵解 $u\in C([0,T];\mathscr{R}_f^+(\Omega))$，它满足相容条件。
\end{theorem}

最后提及下列正则化结果。

\begin{proposition}\label{proposition:ch07-7.3.5}
设 (A0) 与 (A2') 成立。则对几乎所有 $t\in(0,T)$，奇异部分 $u_s(\cdot,t)$ 的支撑集为零测集。
\end{proposition}




\end{document}